\section*{Tag 10 — Strichcode mit Python}

\subsection*{Bleistift-Übung 1: Schachbrett-Regel}

Auf dem Schachbrett ist das Feld an Position $(x, y)$ \textbf{schwarz},
wenn $x + y$ \textbf{gerade} ist, und weiß, wenn $x + y$ ungerade ist.
Mathematisch: $(x + y) \bmod 2 = 0$.

Beweis durch Anschauung: oben links ist $(0, 0)$, $0 + 0 = 0$ — gerade,
also schwarz. Direkt rechts daneben $(1, 0)$, $1 + 0 = 1$ — ungerade,
also weiß. Und direkt darunter $(0, 1)$, $0 + 1 = 1$ — auch weiß. So
wechselt die Farbe diagonal, und das ist die Schachbrett-Regel.

In Python wird daraus die Zeile

\begin{minted}{python}
if (x + y) % 2 == 0:
    brett.putpixel((x, y), 0)
\end{minted}

\subsection*{Bleistift-Übung 2: Hand vs. Python}

Beide Methoden brauchen dasselbe Wissen über das EAN-13-Encoding —
das ist der Punkt: Python nimmt dir das Verstehen nicht ab, nur das
\emph{Tun}. Was sich konkret unterscheidet:

\begin{itemize}
\item \textbf{Zeit pro Code}: Hand etwa 30--60 min für den ersten,
  10--15 min für den zweiten. Python: $< 1$ Sekunde, immer.
\item \textbf{Genauigkeit}: Hand ist auf $\pm 0{,}3$\,mm pro Modul
  begrenzt — die Modulbreite. Python ist pixelgenau.
\item \textbf{Wiederholbarkeit}: Bei Hand sind keine zwei Codes
  exakt gleich. Bei Python ist jedes Bild bis aufs Pixel identisch.
\item \textbf{Lerneffekt}: Hand zeigt dir, was im Strichcode
  \emph{wirklich passiert}. Python zeigt dir, wie du die Idee
  in eine Maschine packst — und wie elegant Code wird, wenn du
  das Problem verstanden hast.
\item \textbf{Skalierung}: Bei einem Code lohnt Python kaum (das
  Skript schreiben dauert länger als das Hand-Zeichnen). Ab dem
  fünften, zehnten, hundertsten Code spart Python alles.
\end{itemize}

Das ist eine Faustregel über Programmieren generell: Code zu schreiben
lohnt sich, sobald du dieselbe Aufgabe mehr als ein paar Mal machst.
Und manchmal lohnt es sich auch nicht — dann ist die Hand schneller.
